\documentclass{article}
\usepackage{amsmath, amssymb, amsthm}
\usepackage{hyperref}
\usepackage{geometry}
\geometry{margin=1in}

\title{Erd\H{o}s Problem \#357}
\date{Last updated: 2025-08-31}
\author{Source: erdosproblems.com}

\begin{document}
\maketitle

\noindent\textbf{Status:} OPEN \quad \textbf{No prize}

\noindent\textbf{Tags:} number theory

\medskip

\noindent\textbf{Problem Statement:}

Let $1\leq a_1&#60;\cdots &#60;a_k\leq n$ be integers such that all sums of the shape $\sum_{u\leq i\leq v}a_i$ are distinct. Let $f(n)$ be the maximal such $k$. How does $f(n)$ grow? Is $f(n)=o(n)$?




Asked by Erd\H{o}s and Harzheim. In \cite{Er77c} Erd\H{o}s asks about an infinite such set of integers, and whether such a set must have density $0$. He notes that a simple averaging process implies $a_k \gg k\log k$ for infinitely many $k$, and so the lower density is $0$. He also asks whether $\sum\frac{1}{a_k}$ must converge. Weisenberg in the comments observes that any set which satisfies [874] also has this property, which implies $f(n)\geq (2+o(1))n^{1/2}$. If $g(n)$ is the maximal $k$ such that there are $1\leq a_1,\ldots,a_k\leq n$ with all consecutive sums distinct (i.e. we drop the monotonicity assumption in the definition of $f$) then Hegyv\'{a}ri \cite{He86} has proved that\[\left(\frac{1}{3}+o(1)\right) n\leq g(n)\leq \left(\frac{2}{3}+o(1)\right)n.\]The upper bound of Coppersmith and Phillips in [867] implies\[g(n) \leq \left(\frac{2}{3}-\frac{1}{512}+o(1)\right)n.\]A similar question can be asked if we replace strict monotonicity with weak monotonicity (i.e. we allow $a_i=a_j$). Erd\H{o}s and Harzheim also ask what is the least $m$ which is not a sum of the given form? Can it be much larger than $n$? Erd\H{o}s and Harzheim can show that $\sum_{x&lt;a_i&lt;x^2}\frac{1}{a_i}\ll 1$. Is it true that $\sum_i \frac{1}{a_i}\ll 1$? See also [34] , [356] , [670] , and [867] . The multiplicative analogue is [421] .




References

[Er77c] Erd\H{o}s, Paul, Problems and results on combinatorial number theory. III . Number theory day (Proc. Conf., Rockefeller Univ., New York, 1976) (1977), 43-72.

[He86] Hegyv\'{a}ri, N., On consecutive sums in sequences . Acta Math. Hungar. (1986), 193--200.

\noindent\textbf{Related OEIS sequences:} A364132, A364153, possible


\bigskip
\noindent\small{Source: \url{https://www.erdosproblems.com/357}}
\end{document}
